\section{Automatic structures: deciding arithmetic with automata}
\label{sec:automatic}

One contributing effort decides first-order statements about unbounded natural
numbers by compiling them to finite automata. The fragment is not all of
arithmetic --- it could not be --- but it is wider than it first appears, and
the certificate is unusually cheap to check.

A predicate is \emph{automatic} when the tuples satisfying it, written in binary
and read synchronously one bit-position at a time, form a regular language.
Linear equalities and inequalities are automatic; so are congruences, powers of
two, popcount parity, and bitwise \textsc{and} and \textsc{xor}. Boolean
connectives become automaton products, and --- the reason the fragment is
interesting --- \emph{quantifiers become projection and determinisation}. An
existential quantifier erases one track and the result is determinised; a
universal is the complement of a projection. Arbitrary quantifier nesting is
therefore supported, at the usual cost.

\subsection{The trailing zeros, and why they are the whole difficulty}

Natural numbers have no canonical bit-width, so the encoding pads. Padding is
where this subject goes wrong, and the effort's own framing of the problem is
worth stating exactly: a witness may need \emph{more digits than any of the
visible inputs}. Projecting away a track without accounting for that loses
witnesses that exist only at greater width.

The repair is \emph{zero-tail saturation}: before determinising a projection,
close the state set under the transitions taken on the all-zeros extension, so
that a state reachable only by reading further padding is still reachable. The
effort's negative control is the smallest possible one --- $y=2x+1$, where the
naive projection is correct for every represented input and wrong at $x=0$ ---
and it is the same shape as the exceptional-cell trap in
Section~\ref{sec:realalgebraic}: correct on everything considered, wrong on the
case dropped by the encoding, and invisible until someone takes a negation.

\subsection{What the checker does, and what it refuses to do}

The certificate is a construction trace: the automaton for each subformula,
together with the local equations relating it to its children --- this product
is that intersection, this projection has these saturated states. The checker
verifies those local equations against \emph{a separately supplied expected
query}, and that is all it does.

It does not explore the automaton, minimise it, or re-run the compiler. The
effort is explicit that its checker ``imports neither producer nor formula
helpers, performs no graph exploration or minimization''. This is the cheapest
producer/checker separation anywhere in the collection: the search does
determinisation, which is exponential; the checker does table lookups.

A second object is the \emph{least-witness graph}. For an existential, the
automaton can be refined so that it computes not merely whether a witness
exists but the least one, as a regular function of the inputs --- which is then
evaluated at concrete naturals by a dynamic program. The recorded corpus does
this at $10^{100}+37$, which is the point: the witness is produced by running a
finite graph, not by searching a bounded range.

\begin{remark}[A regular function is a witness a source proof can use]
This is the same distinction Section~\ref{sec:realalgebraic} draws between a
count and a root index, arrived at from the other side. An automaton that
accepts the graph of a function is a \emph{uniform} witness: one finite object
that answers for every input, rather than a procedure that must be re-run. What
it costs is that the function must be regular, which brings us to the limit.
\end{remark}

\subsection{The limit, proved rather than assumed}

Squaring is not automatic. Neither is exponentiation. The effort proves it by a
growth argument rather than asserting it: a synchronised automaton with a fixed
number of states cannot recognise $y=x^2$, because the relation between the
lengths of the two encodings is not bounded by the pumping structure a fixed
state set provides.

That is the third of the three independent impossibility results this round
produced, and like the other two it says the \emph{representation} is the
obstacle. No amount of additional search inside the automatic fragment reaches
multiplication of two variables. A router should treat a nonlinear atom as a
signal to leave this worker, not as a harder instance of it.

\begin{remark}[Two different reasons a finite object suffices]
It is worth putting this beside Section~\ref{sec:wsts}. There, finiteness comes
from a well-quasi-order: an upward-closed set has a finite basis. Here it comes
from regularity: a language closed under the relevant operations has a finite
automaton. In the same round a third worker gets finiteness from a Bell-number
bound on equality patterns. Three finite certificates for infinite-state
questions, three unrelated reasons, and no one of them is a special case of
another.
\end{remark}

\subsection{Scope}

The fragment is linear arithmetic with the listed bitwise predicates, over
naturals, in a fixed base. Multiplication of variables, and anything that grows
faster than the encoding, are outside it and are outside it provably. Automaton
sizes are the cost: determinising a projection is exponential in the worst case,
and the effort reports its own minimisation ablation on nine examples with the
warning not to learn a cost rule from those nine and then report its improvement
on the same nine.

And as with the rest of this round: no Lean. This effort ships none, and says
that no scaffold is included merely to make the package look implemented.
